% BHU PYQ -- Professional Exam Template
\documentclass[11pt,a4paper]{article}

\usepackage[a4paper,top=2.5cm,bottom=2.5cm,left=2.8cm,right=2.8cm,headheight=14pt]{geometry}
\usepackage{amsmath,amssymb}
\usepackage[T1]{fontenc}
\usepackage[utf8]{inputenc}
\usepackage{lmodern}
\usepackage{microtype}
\usepackage{fancyhdr}
\usepackage{enumitem}
\usepackage{hyperref}
\usepackage{lastpage}
\usepackage{parskip}

\hypersetup{colorlinks=true,urlcolor=black,linkcolor=black,
  pdftitle={B.Sc. (Hons.) Botany Sem-IV BOB-403A 2013-14 -- BHU PYQ}}

\pagestyle{fancy}
\fancyhf{}
\renewcommand{\headrulewidth}{0.4pt}
\renewcommand{\footrulewidth}{0.4pt}
\fancyhead[L]{\small   \textit{Banaras Hindu University}}
\fancyhead[R]{\small   \textit{Botany Paper: BOB-403A (2013-14)}}
\fancyfoot[C]{\small Page  \thepage\ of \pageref{LastPage}}


\newlist{parts}{enumerate}{2}
\setlist[parts,1]{label=(\alph*),leftmargin=*,itemsep=4pt,topsep=3pt}
\setlist[parts,2]{label=(\roman*),leftmargin=*,itemsep=2pt,topsep=2pt}

\newcommand{\pts}[1]{\hfill   \textit{[#1]}}


\begin{document}


\begin{center}
  {\large\bfseries BANARAS HINDU UNIVERSITY}\\[2pt]
  {\normalsize B.Sc. (Hons.) Semester IV Examination, 2013-14}\\[6pt]
  \rule{0.8\linewidth}{0.5pt}\\[6pt]
  {\large\bfseries Botany (Ancillary)}\\[4pt]
  {\large\bfseries Paper: BOB-403A --- Ancillary Botany-II}\\[4pt]
  {\normalsize Time: Three Hours \hfill Maximum Marks: 70}
\end{center}

\rule{\linewidth}{0.4pt}

\smallskip



\noindent   \textit{Note: Attempt five questions in all. Question No. 9 is compulsory. The figures in the right-hand margin indicate marks.}

\bigskip

\medskip


\noindent   \textbf{Question 1.} Briefly describe the disease cycle and control measures of early blight of potato.\pts{15}

\medskip


\noindent   \textbf{Question 2.} Write short notes on any two of the following:\pts{15}

\begin{parts}
  \item   \textit{Ocimum sanctum}
  \item   \textit{Aloe vera}
  \item Papaya mosaic
\end{parts}

\medskip


\noindent   \textbf{Question 3.} Briefly describe photosynthesis.\pts{15}

\medskip


\noindent   \textbf{Question 4.} Briefly describe any two of the following:\pts{15}

\begin{parts}
  \item Glycolysis
  \item Chloroplast
  \item Krebs cycle
\end{parts}

\medskip


\noindent   \textbf{Question 5.} Describe the physiological role of gibberellins in plants.\pts{15}

\medskip


\noindent   \textbf{Question 6.} Write short notes on any two of the following:\pts{15}

\begin{parts}
  \item Food chain
  \item Pyramid of energy
  \item Ecosystem structure
\end{parts}

\medskip


\noindent   \textbf{Question 7.} Describe the various types of interactions in the ecosystem.\pts{15}

\medskip


\noindent   \textbf{Question 8.} Briefly describe any two of the following:\pts{15}

\begin{parts}
  \item Biotic component
  \item Abiotic component
  \item Positive interaction
\end{parts}

\medskip


\noindent   \textbf{Question 9.} Fill in the blanks or answer the following in not more than 50 words each:\pts{10}

\begin{parts}
  \item Citrus canker is caused by the organism \makebox[3cm]{\hrulefill}.
  \item   \textit{Tinospora cordifolia} can be mostly seen on top of \makebox[3cm]{\hrulefill} tree.
  \item Ursolic acid is the chemical constituent of \makebox[3cm]{\hrulefill} medicinal plant.
  \item Calvin cycle takes place in \makebox[3cm]{\hrulefill}.
  \item How many ATPs are produced during glycolysis?
  \item Define commensalism.
  \item Define antibiosis.
  \item What is phototropism?
  \item What are secondary producers?
  \item Define ecosystem.
\end{parts}

\vfill
\rule{\linewidth}{0.4pt}

\begin{center}\small
\href{https://bhu-exam.vercel.app/}{Click Here}\\[4pt]\href{https://me-aryan.vercel.app/}{Click Here}\\[4pt]\href{https://quashan-me.vercel.app/}{Click Here}
\end{center}

\end{document}
